% Options for packages loaded elsewhere
\PassOptionsToPackage{unicode}{hyperref}
\PassOptionsToPackage{hyphens}{url}
\documentclass[
]{article}
\usepackage{xcolor}
\usepackage{amsmath,amssymb}
\setcounter{secnumdepth}{-\maxdimen} % remove section numbering
\usepackage{iftex}
\ifPDFTeX
  \usepackage[T1]{fontenc}
  \usepackage[utf8]{inputenc}
  \usepackage{textcomp} % provide euro and other symbols
\else % if luatex or xetex
  \usepackage{unicode-math} % this also loads fontspec
  \defaultfontfeatures{Scale=MatchLowercase}
  \defaultfontfeatures[\rmfamily]{Ligatures=TeX,Scale=1}
\fi
\usepackage{lmodern}
\ifPDFTeX\else
  % xetex/luatex font selection
\fi
% Use upquote if available, for straight quotes in verbatim environments
\IfFileExists{upquote.sty}{\usepackage{upquote}}{}
\IfFileExists{microtype.sty}{% use microtype if available
  \usepackage[]{microtype}
  \UseMicrotypeSet[protrusion]{basicmath} % disable protrusion for tt fonts
}{}
\makeatletter
\@ifundefined{KOMAClassName}{% if non-KOMA class
  \IfFileExists{parskip.sty}{%
    \usepackage{parskip}
  }{% else
    \setlength{\parindent}{0pt}
    \setlength{\parskip}{6pt plus 2pt minus 1pt}}
}{% if KOMA class
  \KOMAoptions{parskip=half}}
\makeatother
\usepackage{color}
\usepackage{fancyvrb}
\newcommand{\VerbBar}{|}
\newcommand{\VERB}{\Verb[commandchars=\\\{\}]}
\DefineVerbatimEnvironment{Highlighting}{Verbatim}{commandchars=\\\{\}}
% Add ',fontsize=\small' for more characters per line
\newenvironment{Shaded}{}{}
\newcommand{\AlertTok}[1]{\textcolor[rgb]{1.00,0.00,0.00}{\textbf{#1}}}
\newcommand{\AnnotationTok}[1]{\textcolor[rgb]{0.38,0.63,0.69}{\textbf{\textit{#1}}}}
\newcommand{\AttributeTok}[1]{\textcolor[rgb]{0.49,0.56,0.16}{#1}}
\newcommand{\BaseNTok}[1]{\textcolor[rgb]{0.25,0.63,0.44}{#1}}
\newcommand{\BuiltInTok}[1]{\textcolor[rgb]{0.00,0.50,0.00}{#1}}
\newcommand{\CharTok}[1]{\textcolor[rgb]{0.25,0.44,0.63}{#1}}
\newcommand{\CommentTok}[1]{\textcolor[rgb]{0.38,0.63,0.69}{\textit{#1}}}
\newcommand{\CommentVarTok}[1]{\textcolor[rgb]{0.38,0.63,0.69}{\textbf{\textit{#1}}}}
\newcommand{\ConstantTok}[1]{\textcolor[rgb]{0.53,0.00,0.00}{#1}}
\newcommand{\ControlFlowTok}[1]{\textcolor[rgb]{0.00,0.44,0.13}{\textbf{#1}}}
\newcommand{\DataTypeTok}[1]{\textcolor[rgb]{0.56,0.13,0.00}{#1}}
\newcommand{\DecValTok}[1]{\textcolor[rgb]{0.25,0.63,0.44}{#1}}
\newcommand{\DocumentationTok}[1]{\textcolor[rgb]{0.73,0.13,0.13}{\textit{#1}}}
\newcommand{\ErrorTok}[1]{\textcolor[rgb]{1.00,0.00,0.00}{\textbf{#1}}}
\newcommand{\ExtensionTok}[1]{#1}
\newcommand{\FloatTok}[1]{\textcolor[rgb]{0.25,0.63,0.44}{#1}}
\newcommand{\FunctionTok}[1]{\textcolor[rgb]{0.02,0.16,0.49}{#1}}
\newcommand{\ImportTok}[1]{\textcolor[rgb]{0.00,0.50,0.00}{\textbf{#1}}}
\newcommand{\InformationTok}[1]{\textcolor[rgb]{0.38,0.63,0.69}{\textbf{\textit{#1}}}}
\newcommand{\KeywordTok}[1]{\textcolor[rgb]{0.00,0.44,0.13}{\textbf{#1}}}
\newcommand{\NormalTok}[1]{#1}
\newcommand{\OperatorTok}[1]{\textcolor[rgb]{0.40,0.40,0.40}{#1}}
\newcommand{\OtherTok}[1]{\textcolor[rgb]{0.00,0.44,0.13}{#1}}
\newcommand{\PreprocessorTok}[1]{\textcolor[rgb]{0.74,0.48,0.00}{#1}}
\newcommand{\RegionMarkerTok}[1]{#1}
\newcommand{\SpecialCharTok}[1]{\textcolor[rgb]{0.25,0.44,0.63}{#1}}
\newcommand{\SpecialStringTok}[1]{\textcolor[rgb]{0.73,0.40,0.53}{#1}}
\newcommand{\StringTok}[1]{\textcolor[rgb]{0.25,0.44,0.63}{#1}}
\newcommand{\VariableTok}[1]{\textcolor[rgb]{0.10,0.09,0.49}{#1}}
\newcommand{\VerbatimStringTok}[1]{\textcolor[rgb]{0.25,0.44,0.63}{#1}}
\newcommand{\WarningTok}[1]{\textcolor[rgb]{0.38,0.63,0.69}{\textbf{\textit{#1}}}}
\usepackage{longtable,booktabs,array}
\newcounter{none} % for unnumbered tables
\usepackage{calc} % for calculating minipage widths
% Correct order of tables after \paragraph or \subparagraph
\usepackage{etoolbox}
\makeatletter
\patchcmd\longtable{\par}{\if@noskipsec\mbox{}\fi\par}{}{}
\makeatother
% Allow footnotes in longtable head/foot
\IfFileExists{footnotehyper.sty}{\usepackage{footnotehyper}}{\usepackage{footnote}}
\makesavenoteenv{longtable}
\setlength{\emergencystretch}{3em} % prevent overfull lines
\providecommand{\tightlist}{%
  \setlength{\itemsep}{0pt}\setlength{\parskip}{0pt}}
\usepackage{bookmark}
\IfFileExists{xurl.sty}{\usepackage{xurl}}{} % add URL line breaks if available
\urlstyle{same}
\hypersetup{
  pdftitle={World Tubes in Gauged Camera Space: Sublinear Frame Scaling for Dynamic Gaussian Splatting},
  pdfauthor={Anonymous},
  hidelinks,
  pdfcreator={LaTeX via pandoc}}

\title{World Tubes in Gauged Camera Space: Sublinear Frame Scaling for
Dynamic Gaussian Splatting}
\author{Anonymous}
\date{2026-07-22}

\begin{document}
\maketitle

Status: arXiv-style working manuscript with generated LaTeX. The
certified correctness and same-representation scaling tables are
populated, and the three-seed progressive Coffee Martini row is
accepted. The remaining public matrix is paused after the local-memory
incident. This is not a final submission until the pixel-matched/sampler
controls and public-scene breadth in
\texttt{WORLD\_TUBES\_EXPERIMENT\_PLAN.md} are complete.

The frozen public workload now has one 21-row manifest: seven primary
Coffee Martini/control rows, six alternate-triplet rows, six
additional-Neural3D rows, one controlled D-NeRF row, and one separately
labelled deterministic audit. Only the three primary progressive rows
are currently accepted.

\subsection{Abstract}\label{abstract}

Dynamic Gaussian splatting methods render a time-varying scene by
evaluating or deforming primitives at a target timestamp, projecting the
resulting Gaussians, binning them into tiles, resolving visibility, and
compositing the visible contributors. This pipeline is efficient for
isolated frames, but it repeats the same world-side work when a renderer
must produce many temporally nearby samples from a known camera program:
video playback, finite exposure, rolling shutter, dense temporal
supervision, or repeated training/evaluation views.

We introduce \textbf{World Tubes in Gauged Camera Space}, a camera-path
compiler for dynamic Gaussian primitives. The core object is a
sensor-time trace atlas over \texttt{(u,v,t)}: each spacetime primitive
is pulled back through the camera-ray bundle and pushed forward along
the ray fiber to produce a reusable viewport-time footprint, conditional
depth model, support certificate, and adjoint structure. Locally, a
spacetime Gaussian pulled through a camera gauge and integrated over ray
depth yields a UVT footprint by a Schur-complement fiber
marginalization. Globally, camera gauges and event-certified domains
make this construction invariant to depth coordinates on bounded
camera-path chart segments, including the tested orbit segments, finite
exposure, and rolling shutter. The renderer evaluates frames as slices
of the compiled atlas, while training accumulates gradients through a
compiled interval VJP.

On our current STAR UVT / projective interval implementation, the
compiled atlas shows sublinear world-side growth across frame count: the
rerun over \texttt{F=\{4,8,16,32,64,128\}} keeps fixed payload growth at
\texttt{1.0x} while per-frame replay grows \texttt{32.0x} (final payload
ratio \texttt{0.03125}), with final fixed/replay CPU compile, forward,
and backward ratios of \texttt{0.0477}, \texttt{0.181}, and
\texttt{0.392}; trained high-motion traces keep final shared/per-frame
interval-entry ratio below \texttt{0.149}, final trace-count ratio at
\texttt{0.1}, final forward ratio below \texttt{0.266}, and final
backward ratio below \texttt{0.094}. A broad real-video audit covers 10
source-distinct cases, 20 projective-interval trainer payloads,
gradient-preserving compiled-adjoint replacement, and fresh-process
median timing with no-first/projective-total ratios of
\texttt{0.565}/\texttt{0.836}. These results suggest that camera-path
compilation is a practical complement to dynamic Gaussian
representations: it does not claim an information-theoretic sublinear
bound in the number of output pixels, but in the tested training regimes
the dominant bottlenecks scale with trace/event complexity rather than
frame count. The remaining per-pixel shading term is real, but
projection, support, binning, visibility metadata, and backward replay
are amortized across time.

\subsection{1. Introduction}\label{1-introduction}

3D Gaussian Splatting (3DGS) made real-time neural rendering practical
by replacing expensive volumetric ray marching with visibility-aware
anisotropic splatting. Dynamic extensions such as 4D Gaussian Splatting,
Deformable 3D Gaussians, Dynamic 3D Gaussians, and Spacetime Gaussian
Feature Splatting extend the representation through deformation fields,
persistent motion, or spacetime primitives. These methods improve
dynamic scene modeling, but they usually retain a per-target-view
rendering loop: evaluate the primitive state at the requested timestamp,
project to screen, build tile bins, estimate or sort depth order, shade,
composite, and backpropagate through that target render.

This is the wrong computational unit for several common workloads. A
video renderer needs many temporally adjacent frames. A finite-exposure
renderer needs many shutter samples per output image. A rolling-shutter
camera couples image row to capture time. Training often revisits the
same camera path or a small family of nearby paths many times. In these
settings, the output image samples still scale as \texttt{O(F\ H\ W)},
but the world-side work should not have to scale linearly with frame
count \texttt{F}. The paper\textquotesingle s claim is therefore not
that pixels disappear. It is that the expensive training bottlenecks
that dominate dynamic Gaussian pipelines-\/-projection, support, tile
membership, ordering metadata, and backward replay-\/-can scale with
trace/event complexity instead of frame count. In our tested regimes
this produces sublinear measured training-time growth for the compiled
route.

We propose to compile the dynamic world through the camera program
itself. A world primitive is not projected independently into each
frame; it becomes a \textbf{world tube} observed as a
\textbf{sensor-time trace}:

\begin{Shaded}
\begin{Highlighting}[]
\NormalTok{world spacetime primitive {-}\textgreater{} alpha\_i(u,v,t), c\_i(u,v,t), z\_i(u,v,t).}
\end{Highlighting}
\end{Shaded}

The compiled object is a \textbf{trace atlas} over the sensor-time base:

\begin{Shaded}
\begin{Highlighting}[]
\NormalTok{B = Omega x T,      y = (u, v, tau).}
\end{Highlighting}
\end{Shaded}

Each atlas domain stores active primitive traces, support bounds, local
depth models, visibility/order certificates, fallback metadata, and
differentiable state. Frames are slices of this atlas, and
finite-exposure or rolling-shutter images are integrals or row-coupled
samples through the same object.

This construction is not merely a renamed STAR UVT renderer. STAR UVT
supplies the spacetime Gaussian representation and the sparse Metal
execution lineage; the gauged camera-ray formulation supplies the new
compiler semantics. In particular, it pulls primitives through a moving
camera program, marginalizes the ray-depth fiber without discarding
conditional depth, and partitions support-overlap regions at certified
depth-order events. A single raw interval can therefore fail under a
large-motion order crossing while the corresponding
visibility-stratified gauge domains remain replay-equivalent. This
distinction is part of the method and must not be removed by
implementation cleanup.

Our contributions are:

\begin{enumerate}
\def\labelenumi{\arabic{enumi}.}
\item
  \textbf{A camera-ray bundle formulation for dynamic Gaussian
  rendering.} We define a trace as
  \texttt{pi\_*\ Gamma\^{}*\ world\_primitive}, i.e. pull the primitive
  through the camera program and integrate/summarize it along the ray
  fiber.
\item
  \textbf{A local Schur-complement derivation for world tubes.} Under a
  local camera gauge, depth marginalization of a pulled-back spacetime
  Gaussian gives a UVT Gaussian-like footprint and conditional depth
  statistics used for support, visibility, and gradients.
\item
  \textbf{Event-certified gauge domains.} Instead of treating charts as
  ad hoc fitted patches, we use gauge domains that certify projection
  regularity, support validity, depth/order behavior, fallback
  conditions, and backward support.
\item
  \textbf{A lifted visibility gauge atlas.} Footprint traces are pushed
  down to \texttt{(u,v,t)} for support and shading, while visibility is
  compiled from lifted depth/order fields before ray-depth structure is
  discarded. Pairwise support-overlap predicates become sign
  certificates over gauge domains.
\item
  \textbf{A projective interval atlas implementation.} We implement
  interval compressed projective traces with Metal forward and direct
  backward paths for a STAR UVT feature-tube route. Visibility order and
  tile membership are compiled constants during the direct VJP, while
  trace coefficients, opacity, temporal opacity, spatial precision, and
  color remain differentiable.
\item
  \textbf{A sublinear frame-scaling evaluation protocol.} We report
  payload, trace-count, interval-entry, forward, backward, and timing
  ratios against per-frame replay, plus quality/media tethers to verify
  that the compiled route preserves the baseline
  renderer\textquotesingle s output and gradients.
\end{enumerate}

\subsection{2. Related Work}\label{2-related-work}

\textbf{Gaussian splatting.} 3DGS introduced anisotropic 3D Gaussian
primitives and a visibility-aware rasterizer that supports real-time
rendering and optimization. Our work keeps this rasterization motivation
but changes the unit of compilation from one screen at one time to a
sensor-time camera path.

\textbf{Dynamic Gaussian representations.} 4D-GS combines 3D Gaussians
with 4D neural voxels and lightweight deformation prediction. Deformable
3D Gaussians learn a canonical Gaussian scene plus deformation field.
Dynamic 3D Gaussians track persistent Gaussians over time. Spacetime
Gaussian Feature Splatting adds temporal opacity and parametric
motion/rotation to Gaussian primitives. These methods primarily address
the dynamic scene representation. We instead target the repeated
rendering work induced by known camera paths and many temporal samples.

\textbf{Nonlinear cameras, rolling shutter, and ray-space splatting.}
Gaussian Splatting on the Move models blur and rolling shutter under
natural camera motion. 3DGUT replaces the EWA projection approximation
with an unscented transform to support nonlinear cameras, rolling
shutter, and secondary rays. Our method is complementary:
sigma/projective projection helps define or test gauge domains, while
the trace atlas amortizes camera-path work over many samples.

\textbf{Dynamic view-synthesis datasets.} Neural 3D Video, D-NeRF,
HyperNeRF, DyCheck, Technicolor-style light-field videos, and Google
Immersive-style captures provide public dynamic-scene benchmarks. Our
paper should evaluate on public data for comparison, but the central
experiments must also include a controlled synthetic trace suite where
exact ray-fiber integration and visibility events are known.

\subsection{3. Method}\label{3-method}

\subsubsection{3.1 Sensor-time base and camera-ray
bundle}\label{31-sensor-time-base-and-camera-ray-bundle}

Let the sensor-time base be:

\begin{Shaded}
\begin{Highlighting}[]
\NormalTok{B = Omega x T,      y = (u, v, tau).}
\end{Highlighting}
\end{Shaded}

Let the camera program define a ray bundle:

\begin{Shaded}
\begin{Highlighting}[]
\NormalTok{pi: E\_Gamma {-}\textgreater{} B,}
\NormalTok{pi\^{}\{{-}1\}(y) = F\_y.}
\end{Highlighting}
\end{Shaded}

The fiber \texttt{F\_y} is the ray-depth domain over the sensor sample
\texttt{y}. A camera program maps bundle points into world spacetime:

\begin{Shaded}
\begin{Highlighting}[]
\NormalTok{Gamma: E\_Gamma {-}\textgreater{} M,      M = R\^{}3 x R.}
\end{Highlighting}
\end{Shaded}

For a world primitive \texttt{w\_i}, the invariant sensor-time trace is:

\begin{Shaded}
\begin{Highlighting}[]
\NormalTok{Trace\_Gamma[w\_i] = pi\_* Gamma\^{}* w\_i.}
\end{Highlighting}
\end{Shaded}

For a density-like primitive, this becomes:

\begin{Shaded}
\begin{Highlighting}[]
\NormalTok{bar\_rho\_i(y) = integral\_\{F\_y\} rho\_i(Gamma(e)) dmu\_y(e).}
\end{Highlighting}
\end{Shaded}

This equation is the renderer\textquotesingle s invariant object. It
does not depend on a particular depth coordinate or local chart.

\subsubsection{3.2 Gauges and event-certified
domains}\label{32-gauges-and-event-certified-domains}

A gauge is a local trivialization of the ray bundle over a domain
\texttt{C\_a}:

\begin{Shaded}
\begin{Highlighting}[]
\NormalTok{chi\_a: E\_Gamma | C\_a {-}\textgreater{} C\_a x D\_a,}
\NormalTok{chi\_a(e) = (y, z\_a).}
\end{Highlighting}
\end{Shaded}

In gauge \texttt{a}, the trace is:

\begin{Shaded}
\begin{Highlighting}[]
\NormalTok{bar\_rho\_i\^{}a(y)}
\NormalTok{  = integral\_\{D\_a\} rho\_i(Gamma\_a(y, z\_a)) J\_a(y, z\_a) dz\_a.}
\end{Highlighting}
\end{Shaded}

The Jacobian \texttt{J\_a} is not optional. It is the measure correction
that makes ordinary depth, log depth, inverse depth, orbit-angle, and
projective gauges represent the same physical trace. In our current
gauge-invariance artifact, ordinary-depth and log-depth integration of
the same revolving-camera spacetime Gaussian agree to \texttt{3.50e-13}
relative error with the Jacobian; without it, the value error is at
least \texttt{0.600}. The matching gradient artifact agrees to
\texttt{2.33e-12} relative error with the Jacobian; without it, gradient
error is at least \texttt{0.592}.

We use the term \textbf{gauge domain} instead of weak chart. A gauge
domain certifies:

\begin{Shaded}
\begin{Highlighting}[]
\NormalTok{projection denominator regularity}
\NormalTok{trace representation error}
\NormalTok{conservative support bounds}
\NormalTok{tile{-}time active set}
\NormalTok{conditional depth model}
\NormalTok{stable order, commutable order, or fallback}
\NormalTok{interval gates}
\NormalTok{backward support matching forward support}
\end{Highlighting}
\end{Shaded}

For revolving cameras, a projective/orbit gauge such as
\texttt{r\ =\ tan(theta\ /\ 2)} can make traces rational or low-order
over longer intervals than naive frame time. The domain still ends at
real events: denominator zeros, behind-camera transitions, near/far
crossings, support entering/leaving tiles, order swaps, and
disocclusions.

\textbf{Scope of the current implementation.} The compiler and
experiments cover bounded, event-certified orbit segments inside one
regular projective chart. They do not implement chart transitions for
complete \texttt{360°} or repeated \texttt{720°} revolutions. We
therefore make no full-orbit multi-gauge claim in this paper; such a
transition system is future work rather than an untested part of the
method.

\subsubsection{3.3 Local Gaussian fiber
pushforward}\label{33-local-gaussian-fiber-pushforward}

Let a spacetime Gaussian primitive be:

\begin{Shaded}
\begin{Highlighting}[]
\NormalTok{rho\_i(x) = a\_i exp[{-}1/2 (x {-} m\_i)\^{}T Lambda\_i (x {-} m\_i)],}
\NormalTok{x in R\^{}4.}
\end{Highlighting}
\end{Shaded}

In a local gauge, linearize the camera map around \texttt{(y0,\ z0)}:

\begin{Shaded}
\begin{Highlighting}[]
\NormalTok{Gamma\_a(y,z) \textasciitilde{}= x0 + J eta,}
\NormalTok{eta = [delta\_y, delta\_z]\^{}T.}
\end{Highlighting}
\end{Shaded}

Let \texttt{delta\ =\ m\_i\ -\ x0} and
\texttt{g\ =\ J\^{}T\ Lambda\_i\ delta}, partitioned as
\texttt{g\ =\ {[}g\_y,\ g\_z{]}\^{}T}. Then the local exponent is
\texttt{eta\^{}T\ H\ eta\ -\ 2\ g\^{}T\ eta\ +\ delta\^{}T\ Lambda\_i\ delta}.

Plugging the linearized map into the primitive gives a Gaussian over
\texttt{(y,z)} with precision:

\begin{Shaded}
\begin{Highlighting}[]
\NormalTok{H = J\^{}T Lambda\_i J}
\NormalTok{  = [ H\_yy  H\_yz}
\NormalTok{      H\_zy  H\_zz ].}
\end{Highlighting}
\end{Shaded}

Marginalizing the fiber coordinate \texttt{z} yields the UVT precision:

\begin{Shaded}
\begin{Highlighting}[]
\NormalTok{S = H\_yy {-} H\_yz H\_zz\^{}\{{-}1\} H\_zy.}
\end{Highlighting}
\end{Shaded}

The conditional fiber/depth model is:

\begin{Shaded}
\begin{Highlighting}[]
\NormalTok{z\_hat\_i(y) = z0 + H\_zz\^{}\{{-}1\}(g\_z {-} H\_zy (y {-} y0)),}
\NormalTok{Var(z | y) = H\_zz\^{}\{{-}1\}.}
\end{Highlighting}
\end{Shaded}

For a scalar fiber with \texttt{H\_zz\ \textgreater{}\ 0}, an
untruncated local fiber, and locally constant fiber-measure factor
\texttt{J\_0}, the marginal also has amplitude:

\begin{Shaded}
\begin{Highlighting}[]
\NormalTok{bar\_rho\_i(y)}
\NormalTok{  \textasciitilde{}= J\_0 a\_i sqrt(2 pi / H\_zz) exp[{-}1/2 q\_y(delta\_y)],}

\NormalTok{q\_y(delta\_y)}
\NormalTok{  = delta\_y\^{}T S delta\_y}
\NormalTok{    {-} 2 (g\_y {-} H\_yz H\_zz\^{}\{{-}1\} g\_z)\^{}T delta\_y}
\NormalTok{    + delta\^{}T Lambda\_i delta {-} g\_z\^{}T H\_zz\^{}\{{-}1\} g\_z.}
\end{Highlighting}
\end{Shaded}

This local closed form must be replaced by certified quadrature or a
residual bound when fiber clipping or a varying Jacobian is not
negligible.

Thus a pulled-back 4D world Gaussian becomes a 3D sensor-time footprint
plus conditional depth and uncertainty. These quantities are exactly
what a rasterizer needs for support, tile-time binning, visibility
ordering, and gradient propagation.

\subsubsection{3.4 Trace atlas
representation}\label{34-trace-atlas-representation}

The compiled atlas is:

\begin{Shaded}
\begin{Highlighting}[]
\NormalTok{K\_Gamma = \{ C\_l, A\_l, T\_l, Pi\_l, E\_l \}\_\{l=1\}\^{}L.}
\end{Highlighting}
\end{Shaded}

where:

\begin{Shaded}
\begin{Highlighting}[]
\NormalTok{C\_l       gauge domain / event cell in (u,v,t)}
\NormalTok{A\_l       active primitive set}
\NormalTok{T\_l       trace functions: alpha\_i,l(y), c\_i,l(y), z\_i,l(y)}
\NormalTok{Pi\_l      stable total order, partial order, or commutation certificate}
\NormalTok{E\_l       error, support, fallback, and backward metadata}
\end{Highlighting}
\end{Shaded}

Rendering at \texttt{y\ in\ C\_l} evaluates active traces and composites
them in the compiled order:

\begin{Shaded}
\begin{Highlighting}[]
\NormalTok{I(y) = sum\_m T\_m(y) alpha\_\{pi\_m,l\}(y) c\_\{pi\_m,l\}(y).}
\end{Highlighting}
\end{Shaded}

The transmittance is:

\begin{Shaded}
\begin{Highlighting}[]
\NormalTok{T\_m(y) = product\_\{n\textless{}m\} (1 {-} alpha\_\{pi\_n,l\}(y)).}
\end{Highlighting}
\end{Shaded}

A frame is a slice:

\begin{Shaded}
\begin{Highlighting}[]
\NormalTok{I\_k(u,v) = I(u,v,t\_k).}
\end{Highlighting}
\end{Shaded}

Finite exposure is an integral through the atlas:

\begin{Shaded}
\begin{Highlighting}[]
\NormalTok{I\_k(u,v) = integral w\_k(u,v,tau) I(u,v,tau) d tau.}
\end{Highlighting}
\end{Shaded}

Rolling shutter replaces \texttt{tau} with a row/time-coupled sensor
program.

\subsubsection{3.5 Visibility gauge
atlas}\label{35-visibility-gauge-atlas}

Depth marginalization alone does not make alpha compositing linear in
depth. The footprint trace:

\begin{Shaded}
\begin{Highlighting}[]
\NormalTok{pi\_* Gamma\^{}* w\_i}
\end{Highlighting}
\end{Shaded}

is sufficient for support and shading, but visibility is a relation
between lifted traces before the ray-depth structure is fully
marginalized. We therefore compile a second object, a \textbf{visibility
gauge atlas}:

\begin{Shaded}
\begin{Highlighting}[]
\NormalTok{O\_Gamma = \{ C\_l, G\_l, Delta\_l, Pi\_l, R\_l \}.}
\end{Highlighting}
\end{Shaded}

Here \texttt{G\_l} is a local support-overlap graph; it contains only
primitive pairs whose sensor-time footprints overlap inside the same
tile-time cell. \texttt{Delta\_l} stores certified depth/order
predicates, \texttt{Pi\_l} stores the induced total order or
partial-order DAG, and \texttt{R\_l} stores commutation residuals and
fallback metadata. This avoids the all-pairs \texttt{N\^{}2} problem:
pairwise certification is local in support overlap and event complexity,
not global in primitive count.

For each primitive we keep conditional depth:

\begin{Shaded}
\begin{Highlighting}[]
\NormalTok{z\_hat\_i(y),    sigma\_z,i(y).}
\end{Highlighting}
\end{Shaded}

or a conservative lifted interval in a depth/order gauge:

\begin{Shaded}
\begin{Highlighting}[]
\NormalTok{D\_i(y) = [z\_i\^{}{-}(y), z\_i\^{}+(y)].}
\end{Highlighting}
\end{Shaded}

For support-overlapping pairs, define either a center-depth difference:

\begin{Shaded}
\begin{Highlighting}[]
\NormalTok{Delta\_ij(y) = z\_i(y) {-} z\_j(y)}
\end{Highlighting}
\end{Shaded}

or a conservative interval predicate:

\begin{Shaded}
\begin{Highlighting}[]
\NormalTok{Delta\_ij\^{}{-}(y) = z\_i\^{}{-}(y) {-} z\_j\^{}+(y)}
\NormalTok{Delta\_ij\^{}+(y) = z\_i\^{}+(y) {-} z\_j\^{}{-}(y).}
\end{Highlighting}
\end{Shaded}

Then:

\begin{Shaded}
\begin{Highlighting}[]
\NormalTok{if Delta\_ij\^{}+(y) \textless{} 0:  i definitely in front of j}
\NormalTok{if Delta\_ij\^{}{-}(y) \textgreater{} 0:  j definitely in front of i}
\NormalTok{otherwise:             split, commute, or fallback}
\end{Highlighting}
\end{Shaded}

The gauge should be chosen so these depth-difference fields are cheap to
certify. For orbit/projective cameras, ordinary frame time may make
depth curves high-curvature, while projective time, inverse depth, log
depth, or denominator gauges can make them rational, low-degree, or
interval-friendly.

If unresolved pairs remain, we bound the effect of swapping two
translucent contributors:

\begin{Shaded}
\begin{Highlighting}[]
\NormalTok{|Delta I\_ij(y)| \textless{}= alpha\_i(y) alpha\_j(y) |c\_i(y) {-} c\_j(y)|.}
\end{Highlighting}
\end{Shaded}

Unresolved pairs below tolerance are marked commutable. Important
unresolved pairs induce an event boundary or fallback. Fallback is part
of the theorem of the implementation: hard regions can be rendered by
local live sorting or a reference path, while the rest of the atlas
retains shared metadata.

The baseline-compatible theorem is:

\begin{Shaded}
\begin{Highlighting}[]
\NormalTok{Within a domain C\_l, if every support{-}overlapping pair has a certified order}
\NormalTok{or a certified commutation residual below epsilon, then compiled compositing}
\NormalTok{matches the baseline sorted{-}Gaussian renderer up to epsilon and trace}
\NormalTok{approximation error.}
\end{Highlighting}
\end{Shaded}

This claim is intentionally baseline-relative. It reproduces the chosen
Gaussian-splat compositing semantics. It does not claim that
center-depth alpha compositing is a physically exact solution of
radiative transfer.

\subsubsection{3.6 WorldFoam as lifted
transmittance}\label{36-worldfoam-as-lifted-transmittance}

The visibility gauge atlas preserves baseline Gaussian-splat semantics.
A more radical extension is to retain the lifted ray-fiber opacity field
itself:

\begin{Shaded}
\begin{Highlighting}[]
\NormalTok{sigma\_l(y,z) = sum\_i rho\_i(Gamma\_l(y,z)).}
\end{Highlighting}
\end{Shaded}

Instead of sorting primitive centers, one renders by Beer-Lambert
transmittance:

\begin{Shaded}
\begin{Highlighting}[]
\NormalTok{tau(y,z) = integral\_\{z\_front\}\^{}\{z\} sigma\_l(y,s) ds}
\NormalTok{T(y,z)   = exp({-}tau(y,z))}
\NormalTok{I(y)     = integral T(y,z) sigma\_l(y,z) c\_l(y,z) dz.}
\end{Highlighting}
\end{Shaded}

This \textbf{world foam} mode dissolves depth order into cumulative
opacity along the ray fiber. It is cleaner for physical transmittance,
finite exposure, and translucent ambiguity, but it is less directly
baseline-compatible with standard Gaussian splat alpha compositing. We
therefore treat it as a separate representation layer and a second paper
direction: world tubes compile primitive support and differentiable
attributes; WorldFoam compiles lifted opacity/transmittance for
visibility.

\subsubsection{3.7 Compiled adjoints}\label{37-compiled-adjoints}

Inside a gauge domain with fixed support and visibility metadata,
rendering is differentiable with respect to trace parameters. For a
primitive \texttt{i}, the local derivative of compositing is:

\begin{Shaded}
\begin{Highlighting}[]
\NormalTok{dI/dc\_i     = T\_i alpha\_i,}
\NormalTok{dI/dalpha\_i = T\_i (c\_i {-} I\_behind,i).}
\end{Highlighting}
\end{Shaded}

The gradient of the loss is:

\begin{Shaded}
\begin{Highlighting}[]
\NormalTok{dL/dtheta\_i =}
\NormalTok{  sum\_l integral\_\{C\_l\} A\_l(y)\^{}T dI(y)/dtheta\_i dy,}
\end{Highlighting}
\end{Shaded}

where \texttt{A\_l(y)\ =\ dL/dI(y)} is the image adjoint. The compiled
implementation uses interval Metal forward and direct VJP with topology,
active intervals, and visibility cells held as compiled constants. Trace
coefficients, opacity, temporal opacity, spatial precision, and colors
remain differentiable.

\subsection{4. Implementation}\label{4-implementation}

Our current implementation is a STAR UVT / projective interval backend.
A trace stores homogeneous/projective time coefficients, opacity,
optional temporal opacity coefficients, optional spatial precision,
optional depth-affine terms, and color. Tile-time cells store active
intervals and visibility metadata. The hot path packs accepted cells
once into spatial tile bins and uses per-entry
\texttt{{[}active\_start,\ active\_stop)} checks in the Metal kernel.

The forward path is:

\begin{Shaded}
\begin{Highlighting}[]
\NormalTok{render\_projective\_trace\_cell\_interval\_atlas\_metal(...)}
\end{Highlighting}
\end{Shaded}

The direct VJP path is:

\begin{Shaded}
\begin{Highlighting}[]
\NormalTok{direct\_backward\_projective\_trace\_cell\_interval\_atlas\_metal(...)}
\end{Highlighting}
\end{Shaded}

The trainer route calls:

\begin{Shaded}
\begin{Highlighting}[]
\NormalTok{\_render\_projective\_interval\_feature\_tubes\_autograd(...)}
\end{Highlighting}
\end{Shaded}

via a custom autograd function:

\begin{Shaded}
\begin{Highlighting}[]
\NormalTok{\_ProjectiveCellIntervalBackward}
\end{Highlighting}
\end{Shaded}

Current limitations of this implementation:

\begin{Shaded}
\begin{Highlighting}[]
\NormalTok{MPS/Metal backend}
\NormalTok{RGB / feature\_dim=3 route for the projective interval trainer}
\NormalTok{compiled visibility/order and tile membership held fixed during direct VJP}
\NormalTok{fallback support present, but broad fallback{-}heavy scenes remain a stress case}
\NormalTok{STAR UVT support/visibility/composition quality is a separate active research lane}
\end{Highlighting}
\end{Shaded}

These are implementation limits, not limits of the bundle formulation.

\subsection{5. Experiments}\label{5-experiments}

The final paper should separate four questions:

\begin{enumerate}
\def\labelenumi{\arabic{enumi}.}
\tightlist
\item
  \textbf{Correctness:} Does the trace integral match dense ray/fiber
  evaluation?
\item
  \textbf{Sublinear scaling:} Does world-side work grow sublinearly with
  frame count or camera-family size?
\item
  \textbf{Renderer equivalence:} Does the compiled route preserve
  images, losses, media, and gradients relative to the baseline
  renderer?
\item
  \textbf{Usefulness:} Does the speed/memory benefit hold on public
  dynamic-scene workloads without excessive fallback?
\end{enumerate}

\subsubsection{5.1 Synthetic trace
correctness}\label{51-synthetic-trace-correctness}

Controlled scenes of analytic 4D Gaussians with known camera programs:

\begin{Shaded}
\begin{Highlighting}[]
\NormalTok{camera paths: static, dolly, orbit, fast orbit, rolling shutter}
\NormalTok{FOV: 30, 60, 90, 120 degrees}
\NormalTok{frames/samples: 4, 8, 16, 32, 64, 128}
\NormalTok{primitive anisotropy: isotropic, elongated spatial, elongated temporal}
\NormalTok{visibility: isolated, crossing pair, thin foreground occluder, dense translucent}
\end{Highlighting}
\end{Shaded}

Metrics:

\begin{Shaded}
\begin{Highlighting}[]
\NormalTok{trace L1 / KL against dense fiber integration}
\NormalTok{image PSNR / SSIM / LPIPS}
\NormalTok{alpha error}
\NormalTok{depth{-}order error}
\NormalTok{fallback fraction}
\NormalTok{support over/under{-}coverage}
\NormalTok{gradient relative error}
\end{Highlighting}
\end{Shaded}

\subsubsection{5.2 Frame-count scaling}\label{52-frame-count-scaling}

Compare per-frame replay against compiled atlas for fixed camera paths:

\begin{Shaded}
\begin{Highlighting}[]
\NormalTok{F = 4, 8, 16, 32, 64, 128}
\end{Highlighting}
\end{Shaded}

Metrics:

\begin{Shaded}
\begin{Highlighting}[]
\NormalTok{payload bytes}
\NormalTok{trace count}
\NormalTok{tile/bin entries}
\NormalTok{interval entries}
\NormalTok{CPU compile time}
\NormalTok{GPU forward time}
\NormalTok{GPU backward time}
\NormalTok{total step time}
\NormalTok{peak memory}
\end{Highlighting}
\end{Shaded}

Current internal evidence:

\begin{Shaded}
\begin{Highlighting}[]
\NormalTok{bounded{-}orbit F: 4, 8, 16, 32, 64, 128}
\NormalTok{fixed payload growth: 1.0x vs per{-}frame replay 32.0x}
\NormalTok{final fixed/replay payload ratio: 0.03125}
\NormalTok{final fixed/replay CPU compile ratio: 0.0477}
\NormalTok{final fixed/replay forward ratio: 0.181}
\NormalTok{final fixed/replay backward ratio: 0.392}
\NormalTok{trained interval{-}entry growth ratio: 0.148}
\NormalTok{final trained trace{-}count ratio: 0.1}
\NormalTok{final trained forward ratio: \textless{}= 0.266}
\NormalTok{final trained backward ratio: \textless{}= 0.094}
\NormalTok{fresh{-}process median no{-}first ratio: 0.565}
\NormalTok{fresh{-}process median projective{-}total ratio: 0.836}
\end{Highlighting}
\end{Shaded}

\subsubsection{5.3 Camera-family
scaling}\label{53-camera-family-scaling}

For a low-dimensional camera family \texttt{Q\ x\ Omega\ x\ T}, compare
one shared family atlas against replaying one atlas per camera parameter
sample:

\begin{Shaded}
\begin{Highlighting}[]
\NormalTok{Q1: orbit phase offset}
\NormalTok{Q2: orbit phase + height}
\NormalTok{grid: 3x3, 5x5, 7x7}
\end{Highlighting}
\end{Shaded}

Current internal evidence:

\begin{Shaded}
\begin{Highlighting}[]
\NormalTok{Q2 shared payload growth: 1.0x}
\NormalTok{Q2 replay payload growth: 64.0x}
\NormalTok{Q2 final payload ratio: 0.0625}
\NormalTok{Q2 final chart ratio: 0.015625}
\NormalTok{Q2 max UV fit residual: 0.111 px}
\end{Highlighting}
\end{Shaded}

\subsubsection{5.4 Real-video renderer
equivalence}\label{54-real-video-renderer-equivalence}

Use source-distinct real videos and compare compiled projective interval
renderer against the cadence/per-frame route.

Metrics:

\begin{Shaded}
\begin{Highlighting}[]
\NormalTok{loss curve delta}
\NormalTok{RGB loss curve delta}
\NormalTok{end PSNR delta}
\NormalTok{contact{-}sheet pixel delta}
\NormalTok{PNG hash match}
\NormalTok{gradient flags present}
\NormalTok{support rebins}
\NormalTok{stale refreshes}
\NormalTok{fallback fraction}
\end{Highlighting}
\end{Shaded}

Current internal evidence:

\begin{Shaded}
\begin{Highlighting}[]
\NormalTok{10 source{-}distinct broad quality pairs}
\NormalTok{10 source{-}distinct broad media pairs}
\NormalTok{20 broad10 trainer case payloads}
\NormalTok{all cases projective interval main path}
\NormalTok{all renderer gradient flags present}
\NormalTok{zero support rebins / stale refreshes in accepted rows}
\NormalTok{compiled trainer replacement gap: 0}
\end{Highlighting}
\end{Shaded}

\subsubsection{5.5 Public dataset
comparison}\label{55-public-dataset-comparison}

The paper should evaluate public data in two ways:

\begin{enumerate}
\def\labelenumi{\arabic{enumi}.}
\item
  \textbf{In-representation ablation:} train or initialize our STAR
  UVT/projective primitives on public sequences, then compare
  frame-by-frame replay versus compiled atlas. This is the cleanest
  evaluation of the contribution.
\item
  \textbf{External dynamic-GS comparison:} compare speed/quality against
  existing dynamic Gaussian baselines where feasible, but present this
  as contextual comparison rather than the main theorem.
\end{enumerate}

Recommended datasets:

\begin{Shaded}
\begin{Highlighting}[]
\NormalTok{Neural 3D Video: real multiview dynamic scenes}
\NormalTok{D{-}NeRF: controlled synthetic dynamic scenes}
\NormalTok{HyperNeRF / DyCheck: monocular/topology stress tests}
\NormalTok{Technicolor{-}style light{-}field scenes: synchronized camera array stress}
\NormalTok{internal broad10 real{-}video set: engineering{-}only unless made reproducible}
\end{Highlighting}
\end{Shaded}

\subsubsection{5.6 Finite exposure and rolling
shutter}\label{56-finite-exposure-and-rolling-shutter}

Compare:

\begin{Shaded}
\begin{Highlighting}[]
\NormalTok{baseline per{-}shutter{-}sample rendering}
\NormalTok{compiled atlas sampled at same shutter samples}
\NormalTok{compiled atlas with adaptive temporal quadrature}
\NormalTok{high{-}sample reference}
\end{Highlighting}
\end{Shaded}

Metrics:

\begin{Shaded}
\begin{Highlighting}[]
\NormalTok{quality vs high{-}sample reference}
\NormalTok{unique time samples}
\NormalTok{payload growth}
\NormalTok{forward/backward time}
\NormalTok{rolling row{-}time correctness}
\end{Highlighting}
\end{Shaded}

\subsubsection{5.7 Visibility stress
tests}\label{57-visibility-stress-tests}

Create scenes where the compiler should fail locally:

\begin{Shaded}
\begin{Highlighting}[]
\NormalTok{crossing translucent slabs}
\NormalTok{thin foreground occluders}
\NormalTok{near{-}camera large splats}
\NormalTok{fast orbit around high{-}depth{-}variance geometry}
\NormalTok{disocclusion boundary}
\end{Highlighting}
\end{Shaded}

Metrics:

\begin{Shaded}
\begin{Highlighting}[]
\NormalTok{fallback fraction}
\NormalTok{order{-}strata count}
\NormalTok{commutation{-}bound accepted pairs}
\NormalTok{image error near event boundaries}
\NormalTok{speed lost to fallback}
\end{Highlighting}
\end{Shaded}

\subsection{6. Tables and figures}\label{6-tables-and-figures}

Planned paper figures:

\begin{enumerate}
\def\labelenumi{\arabic{enumi}.}
\tightlist
\item
  \textbf{Concept figure:} per-frame dynamic GS replay vs world-tube
  trace atlas.
\item
  \textbf{Bundle diagram:} \texttt{B\ =\ Omega\ x\ T}, ray fibers,
  \texttt{Gamma}, pullback, pushforward.
\item
  \textbf{Schur complement diagram:} 4D primitive -\textgreater{}
  \texttt{(u,v,t,z)} Gaussian -\textgreater{} depth-marginalized UVT
  footprint.
\item
  \textbf{Gauge-domain/event diagram:} orbit camera with projective
  domains split by denominator/support/order events.
\item
  \textbf{System diagram:} compile, atlas, Metal interval forward,
  direct VJP.
\item
  \textbf{Scaling chart:} frame count vs payload/bin/forward/backward
  ratio.
\item
  \textbf{Camera-family chart:} Q-grid replay vs shared family atlas.
\item
  \textbf{Quality tether chart:} compiled vs baseline loss/PSNR/media
  deltas.
\item
  \textbf{Fallback stress chart:} fallback fraction vs visibility
  density.
\end{enumerate}

Core tables:

\begin{Shaded}
\begin{Highlighting}[]
\NormalTok{Table 1: synthetic correctness vs dense fiber reference}
\NormalTok{Table 2: frame{-}count scaling, internal + public sequences}
\NormalTok{Table 3: ablation of trace representation and gauge domains}
\NormalTok{Table 4: training/backward replacement and gradient errors}
\NormalTok{Table 5: public dynamic{-}scene comparison}
\NormalTok{Table 6: finite exposure / rolling shutter}
\NormalTok{Table 7: limitations and fallback{-}heavy cases}
\end{Highlighting}
\end{Shaded}

\subsubsection{6.1 Certified correctness and theorem
table}\label{61-certified-correctness-and-theorem-table}

All rows below are generated from verifier-accepted JSON artifacts.
Their scope is bounded event-certified projective chart segments; they
do not assert an unimplemented full \texttt{360/720} multi-chart
transition.

{\def\LTcaptype{none} % do not increment counter
\begin{longtable}[]{@{}llrr@{}}
\toprule\noalign{}
Claim & Metric & Value & Acceptance \\
\midrule\noalign{}
\endhead
\bottomrule\noalign{}
\endlastfoot
Fiber value is gauge invariant & max relative error &
\texttt{3.50087e-13} & \texttt{\textless{}=\ 1e-10} \\
Fiber gradient is gauge invariant & max gradient relative error &
\texttt{2.32523e-12} & \texttt{\textless{}=\ 1e-9} \\
Compiled atlas matches dense/replay image & max absolute image error &
\texttt{0} & \texttt{\textless{}=\ 1e-5} \\
Unstratified interval exposes an order-crossing failure & raw crossing
quality error & \texttt{0.186742} & \texttt{\textgreater{}\ 1e-5}
(expected failure) \\
Visibility crossing is repaired by stratification & stratified crossing
quality error & \texttt{0} & \texttt{\textless{}=\ 1e-5} \\
Finite exposure / rolling shutter forward parity & max Metal absolute
error & \texttt{5.96046e-08} & \texttt{\textless{}=\ 1e-5} \\
Finite exposure / rolling shutter gradient parity & max Metal gradient
relative error & \texttt{6.37738e-07} & \texttt{\textless{}=\ 1e-5} \\
Mixed fallback preserves gradients & max mixed gradient relative error &
\texttt{7.40632e-07} & \texttt{\textless{}=\ 1e-5} \\
Bounded-orbit chart reuses payload at \texttt{F=128} & fixed/replay
trace ratio & \texttt{0.03125} & \texttt{\textless{}\ 0.25} \\
Bounded-orbit compiled forward is faster at \texttt{F=128} &
fixed/replay forward ratio & \texttt{0.181323} &
\texttt{\textless{}\ 0.5} \\
Bounded-orbit compiled backward is faster at \texttt{F=128} &
fixed/replay backward ratio & \texttt{0.392235} &
\texttt{\textless{}\ 0.5} \\
\end{longtable}
}

\subsubsection{6.2 Exact same-representation frame
scaling}\label{62-exact-same-representation-frame-scaling}

The accepted \texttt{F=\{4,8,16,32,64,128\}} experiment compares
per-frame STAR replay with the compiled projective atlas at identical
representation settings. Fixed payload growth is \texttt{1.0x} while
replay grows \texttt{32.0x}; at \texttt{F=128}, fixed/replay payload,
compile, forward, and backward ratios are \texttt{0.03125},
\texttt{0.047677}, \texttt{0.181323}, and \texttt{0.392235}. This is the
central causal systems result. Public quality rows test whether that
compiler result survives real scene breadth; they are not substitutes
for this same-representation comparison.

\subsubsection{6.3 Public comparison
status}\label{63-public-comparison-status}

The shared progressive/fixed/global-shuffle Coffee Martini protocols,
evidence schema, and matrix generator are complete. Three clean-source
progressive-512 runs (seeds 17, 29, and 43) completed for all
representations on the full 300-frame sequence. The table reports final
checkpoints and equal optimizer, target-frame, and target-pixel budgets.
Storage and parameter counts are not matched: World Tubes shares
temporal trace state, while dynamic 3DGS and WorldFoam retain
substantial per-frame state.

{\def\LTcaptype{none} % do not increment counter
\begin{longtable}[]{@{}lrrrrrrr@{}}
\toprule\noalign{}
Representation & Heldout PSNR & SSIM & LPIPS & L1 & Train wall (s) &
Peak driver (GB) & Checkpoint (MB) \\
\midrule\noalign{}
\endhead
\bottomrule\noalign{}
\endlastfoot
World Tubes & \texttt{5.9153\ +/-\ 0.0053} & \texttt{0.03549} &
\texttt{0.98305} & \texttt{0.45120} & \texttt{78.33\ +/-\ 21.00} &
\texttt{3.114} & \texttt{0.060} \\
WorldFoam & \texttt{5.6159\ +/-\ 0.0083} & \texttt{0.00460} &
\texttt{0.97054} & \texttt{0.47221} & \texttt{361.82\ +/-\ 55.90} &
\texttt{15.794} & \texttt{116.748} \\
Dynamic 3DGS & \texttt{4.9110\ +/-\ 0.0001} & \texttt{0.28267} &
\texttt{0.90228} & \texttt{0.52139} & \texttt{79.44\ +/-\ 6.59} &
\texttt{20.557} & \texttt{17.206} \\
\end{longtable}
}

These are accepted evidence rows, but not a complete public comparison.
The absolute reconstruction quality is low and the metrics disagree:
World Tubes has the best PSNR/L1, whereas dynamic 3DGS has the best
SSIM/LPIPS. The pixel-matched fixed control, global-shuffle control,
additional camera triplets, two additional Neural3D scenes, and
controlled D-NeRF row remain mandatory before submission claims are
frozen.

The next fixed-512 run was killed after severe unified-memory
compression and swap pressure destabilized the local workstation. Its
partial outputs are excluded. The runner now isolates representations in
child processes and is fail-closed on local MPS, but full-scale local
execution remains unauthorized until targets/rays/evaluation are
streamed or a larger machine is used.

\subsection{7. Discussion and
limitations}\label{7-discussion-and-limitations}

The method is not a universal replacement for dynamic Gaussian scene
representations. It is a renderer/compiler layer for known or
low-dimensional camera programs. It is most useful when many temporal
samples reuse a smooth camera path or path family.

It is also not an information-theoretic claim that total work is
sublinear in materialized output pixels. The sharper claim is empirical
and architectural: in regimes where dynamic Gaussian training is
dominated by world-side projection, support, binning, visibility, and
backward replay, compiling those events into world tubes makes the
dominant bottlenecks scale with trace/event complexity instead of frame
count. In our tested regime this yields sublinear end-to-end
training-time growth, while preserving a residual per-pixel shading
term.

Failure modes:

\begin{Shaded}
\begin{Highlighting}[]
\NormalTok{fallback{-}heavy visibility chaos can erase speedups}
\NormalTok{very wide FOV / near{-}camera splats require more gauge domains}
\NormalTok{single random novel views do not amortize compile cost}
\NormalTok{current implementation is RGB/MPS/STAR{-}UVT scoped}
\NormalTok{current broad real{-}video evidence proves renderer equivalence, not general SOTA quality}
\NormalTok{world foam may be a better alpha/transmittance model, but it changes semantics}
\end{Highlighting}
\end{Shaded}

The main research claim is narrower and stronger:

\begin{Shaded}
\begin{Highlighting}[]
\NormalTok{known camera programs expose repeated world{-}side rendering work;}
\NormalTok{camera{-}gauged world tubes compile that work into reusable sensor{-}time traces.}
\end{Highlighting}
\end{Shaded}

The hierarchy is:

\begin{Shaded}
\begin{Highlighting}[]
\NormalTok{World Tubes:}
\NormalTok{  sublinear camera{-}path compilation for dynamic Gaussian{-}splat semantics.}

\NormalTok{Visibility Gauge Atlas:}
\NormalTok{  certified depth/order compilation for baseline{-}compatible alpha compositing.}

\NormalTok{World Foam:}
\NormalTok{  lifted opacity/transmittance compilation that avoids discrete depth sorting}
\NormalTok{  and moves toward volumetric ray{-}fiber transport.}
\end{Highlighting}
\end{Shaded}

\subsection{8. Conclusion}\label{8-conclusion}

We presented World Tubes in Gauged Camera Space, a trace-atlas renderer
for dynamic Gaussian splatting. By defining rendering as a camera-ray
bundle pushforward, deriving local Schur-complement UVT footprints, and
compiling event-certified gauge domains into interval tile-time
metadata, the method shares projection, support, binning, visibility,
and backward work across time. The current implementation demonstrates
sublinear world-side scaling and a compiled-adjoint training route on
STAR UVT/projective interval traces. The next step toward publication is
a public benchmark suite that reproduces these scaling claims on
standard dynamic-view datasets and controlled visibility stress tests.

\subsection{References To Cite}\label{references-to-cite}

\begin{itemize}
\tightlist
\item
  Kerbl et al., "3D Gaussian Splatting for Real-Time Radiance Field
  Rendering," SIGGRAPH 2023. \url{https://arxiv.org/abs/2308.04079}
\item
  Wu et al., "4D Gaussian Splatting for Real-Time Dynamic Scene
  Rendering," CVPR 2024. \url{https://arxiv.org/abs/2310.08528}
\item
  Yang et al., "Deformable 3D Gaussians for High-Fidelity Monocular
  Dynamic Scene Reconstruction," 2023.
  \url{https://arxiv.org/abs/2309.13101}
\item
  Luiten et al., "Dynamic 3D Gaussians: Tracking by Persistent Dynamic
  View Synthesis," 2023. \url{https://arxiv.org/abs/2308.09713}
\item
  Li et al., "Spacetime Gaussian Feature Splatting for Real-Time Dynamic
  View Synthesis," CVPR 2024. \url{https://arxiv.org/abs/2312.16812}
\item
  Jiang et al., "Gaussian Splatting on the Move: Blur and Rolling
  Shutter Compensation for Natural Camera Motion," 2024.
  \url{https://arxiv.org/abs/2403.13327}
\item
  Wu et al., "3DGUT: Enabling Distorted Cameras and Secondary Rays in
  Gaussian Splatting," CVPR 2025. \url{https://arxiv.org/abs/2412.12507}
\item
  Li et al., "Neural 3D Video Synthesis from Multi-view Video," CVPR
  2022. \url{https://arxiv.org/abs/2103.02597}
\item
  Pumarola et al., "D-NeRF: Neural Radiance Fields for Dynamic Scenes,"
  CVPR 2021. \url{https://arxiv.org/abs/2011.13961}
\item
  Park et al., "HyperNeRF: A Higher-Dimensional Representation for
  Topologically Varying Neural Radiance Fields," SIGGRAPH Asia 2021.
  \url{https://hypernerf.github.io/}
\item
  Hou et al., "Sort-free Gaussian Splatting via Weighted Sum Rendering,"
  ICLR 2025. \url{https://arxiv.org/abs/2410.18931}
\item
  Koo et al., "Gaussian Blending: Rethinking Alpha Blending in 3D
  Gaussian Splatting," AAAI 2026.
  \url{https://doi.org/10.1609/aaai.v40i7.37495}
\end{itemize}

\end{document}
